\section{Détection : filtres adaptés et test du rapport de vraisemblance}\label{sec-detection}

\subsection{Gabarits}
Chaque mécanisme de la \cref{sec-anatomy} laisse une forme connue dans l'espace de portage. Le détecteur teste un
\emph{gabarit} $t\in\R^K$ par mécanisme et par position :
\begin{center}\small
\begin{tabular}{lp{3.3cm}p{4.3cm}p{4.6cm}}
\toprule
Gabarit & Vecteur & Mécanismes & Amplitude $a$ en unités naturelles\\
\midrule
niveau & $\one$ & décalage du spot ; dividende avant $T_1$ & erreur de spot $\approx-a$ ; dividende proche $\approx-a(S-D_1)$\\
marche $k$ & $\one\{T\ge T_k\}$ & montant de dividende dans $(T_{k-1},T_k]$ ; changement de panel en $T_k$ & erreur de dividende $\approx-a(S-D_k)$\\
pic $k$ & $e_k$ & forward aberrant ; date de détachement de part et d'autre de $T_k$ & $\delta\ln F_k\approx-a$ ; part du dividende $\approx-a(S-D_k)/\tilde d_i$\\
saison & $\sum_i\tilde d_iK_h(T_k-t_i)/(S-D_k)$ & décalage du profil de détachement & décalage $\approx365\,a$ jours\\
\bottomrule
\end{tabular}
\end{center}
$K_h$ est un noyau gaussien ($h=20$ jours). Les unités naturelles découlent des
\cref{prop-spot,prop-bucket,prop-timing,prop-profile}. Un pic au dernier pilier est le même vecteur qu'une marche à ce
pilier et il est présenté comme l'un ou l'autre. Avec onze piliers, on obtient \vNTemplates{} gabarits.

\subsection{Blanchiment, projection et filtre adapté}
Soit $\kappa\Sigma=LL^\top$ (Cholesky), $\tilde y=L^{-1}y$, $\tilde X=L^{-1}X$, et $P=I-QQ^\top$ avec $Q$ une base
orthonormée des colonnes de $\tilde X$ : $P$ retire le modèle nul dans l'espace blanchi. On note $r=P\tilde y$ et
$\tb=PL^{-1}t$.

\begin{theorem}[Test du rapport de vraisemblance généralisé pour un gabarit]\label{thm-glrt}
Dans le modèle $y=X\beta+a\,t+\varepsilon$, $\varepsilon\sim\N(0,\kappa\Sigma)$, avec $\beta$ inconnu, la statistique du
rapport de vraisemblance généralisé de $H_0:a=0$ contre $H_1:a\ne0$ est $z_t^2$, avec
\begin{equation}\label{eq-z}
z_t=\frac{\tb^\top r}{\|\tb\|},\qquad z_t\sim\N\bigl(a\|\tb\|,\,1\bigr).
\end{equation}
L'amplitude du maximum de vraisemblance est $\hat a_t=\tb^\top r/\|\tb\|^2$, d'écart-type $1/\|\tb\|$. Parmi les tests
invariants par le sous-espace du modèle nul et par le signe de $a$, le test qui rejette pour les grandes valeurs de
$|z_t|$ est uniformément le plus puissant.
\end{theorem}
\begin{proof}
Après blanchiment, $\tilde y=\tilde X\beta+a\,L^{-1}t+\tilde\varepsilon$ avec $\tilde\varepsilon\sim\N(0,I)$. La
log-vraisemblance maximisée vaut $-\frac12$ fois la somme des carrés des résidus, soit $\|P\tilde y\|^2$ sous $H_0$ et le
résidu de la régression sur $[\tilde X,L^{-1}t]$ sous $H_1$. D'après le théorème de Frisch--Waugh--Lovell~\cite{lovell1963},
ce dernier vaut $\|r\|^2-(\tb^\top r)^2/\|\tb\|^2$, donc $2\ln\Lambda=(\tb^\top r)^2/\|\tb\|^2=z_t^2$, et $\hat a_t$ est le
coefficient de la régression partielle. Comme $P\tilde X=0$ et $\tb^\top P=\tb^\top$,
$\tb^\top r=a\,\tb^\top L^{-1}t+\tb^\top\tilde\varepsilon=a\|\tb\|^2+\tb^\top\tilde\varepsilon$, et
$\tb^\top\tilde\varepsilon/\|\tb\|\sim\N(0,1)$. L'optimalité est le résultat des détecteurs de sous-espace adapté de
Scharf et Friedlander~\cite{scharf1994} ; voir aussi~\cite{kay1998}.
\end{proof}

\begin{corollary}[Puissance et rapport signal sur bruit]\label{cor-power}
Au seuil $c$, la probabilité de détection vaut $\Phi(\lambda-c)+\Phi(-\lambda-c)$, où $\lambda=|a|\,\|\tb\|$ est le rapport
signal sur bruit de l'anomalie. Avec le seuil familial typique du banc d'essai, $c\approx\vZcrit$, une anomalie de
$\lambda=4$ est détectée avec une probabilité de \vPowThFour{} et une anomalie de $\lambda=6$ avec \vPowThSix.
\end{corollary}
\begin{proof}
$|z_t|>c$ avec $z_t\sim\N(\lambda,1)$ (au signe près).
\end{proof}

\begin{proposition}[Test global]\label{prop-q}
Sous $H_0$, $Q=\|r\|^2\sim\chi^2_{K-p}$, où $p=2$ est le nombre de colonnes du modèle nul.
\end{proposition}
\begin{proof}
$r=P\tilde\varepsilon$, où $P$ est une projection orthogonale de rang $K-p$.
\end{proof}
$Q$ détecte une rugosité qu'aucun gabarit n'explique.

\subsection{Calibration familiale}
\begin{proposition}[Valeur critique familiale]\label{prop-family}
Sous $H_0$, le vecteur $z=(z_t)_t$ est gaussien centré de corrélation $C_{tu}=\tb^\top\tilde u/(\|\tb\|\|\tilde u\|)$, qui ne
dépend ni de $\beta$ ni de $\kappa$. La valeur critique $c_\alpha$ telle que $P(\max_t|z_t|>c_\alpha)=\alpha/2$ se calcule
exactement en simulant $\N(0,C)$ ; elle est au plus égale à la valeur de Šidák
$\Phi^{-1}\bigl(1-\tfrac12[1-(1-\alpha/2)^{1/m}]\bigr)$ pour $m$ gabarits. La règle « alerte si $\max_t|z_t|>c_\alpha$ ou si
$Q$ dépasse son quantile $1-\alpha/2$ » a un taux de fausses alertes d'au plus $\alpha$.
\end{proposition}
\begin{proof}
$z=D^{-1}\tilde T^\top P\tilde\varepsilon$, où $\tilde T$ est la matrice des $\tb$ et $D$ celle de leurs normes, d'où la
corrélation. L'inégalité de Šidák~\cite{sidak1967} donne $P(\max_t|z_t|\le c)\ge\prod_tP(|z_t|\le c)$ pour toute
corrélation : le seuil de Šidák est donc conservateur. Les deux tests se partagent le budget $\alpha$ par l'inégalité de
la réunion.
\end{proof}

\begin{proposition}[Calibration exacte de l'échelle de variance]\label{prop-kappa}
Remplacer $\kappa=1$ par $\kappa$ divise chaque $z_t$ par $\sqrt\kappa$ et $Q$ par $\kappa$, et laisse $c_\alpha$ inchangé.
Pour un ensemble de sous-jacents sains, l'indicatrice d'alerte est donc décroissante en $\kappa$, le taux d'alerte
empirique est une fonction en escalier décroissante de $\kappa$, et le plus petit $\kappa$ dont le taux d'alerte ne dépasse
pas la cible s'obtient exactement par dichotomie sur $\ln\kappa$.
\end{proposition}
\begin{proof}
Blanchir avec $\sqrt\kappa L$ multiplie $\tilde y$, $\tilde X$ et chaque $\tb$ par $\kappa^{-1/2}$ ; $P$ est inchangé, donc $r$ et
$\tb$ sont multipliés par $\kappa^{-1/2}$, $z_t$ par $\kappa^{-1/2}$ et $Q$ par $\kappa^{-1}$ ; $C$ est une corrélation et ne bouge
pas. L'indicatrice $\{\max|z|/\sqrt\kappa>c\}\cup\{Q/\kappa>q\}$ de chaque sous-jacent se réduit quand $\kappa$ augmente.
\end{proof}

\subsection{Attribution, réparation et test d'erreur de marquage FO}
Après une alerte, les gabarits sont retenus de façon gloutonne (poursuite adaptative~\cite{mallat1993}) : à chaque
étape, on ajoute le gabarit de plus grand $|z|$ après projection du modèle nul et des gabarits déjà retenus, tant qu'il
dépasse $c_\alpha$. Les gabarits retenus sont ensuite réajustés conjointement avec le modèle nul par moindres carrés
généralisés, chaque amplitude est traduite en unités naturelles, et le portage \emph{réparé} vaut
$y-\sum_{t\in\mathcal S}\hat a_t\,t$. Les gabarits contenus dans le sous-espace du modèle nul sont indétectables par
construction : un écart de repo lisse est un désaccord légitime, publié mais jamais signalé comme anomalie.

Le niveau de repo $\hat\beta_0$ (Totem moins FO à un an) et son écart-type proviennent du même ajustement. Le repo FO est
signalé comme erreur de marquage possible lorsque $|\hat\beta_0|$ dépasse à la fois $10$\,pb et
$t_{\nu,1-\alpha/2}\cdot s\cdot\operatorname{se}(\hat\beta_0)$ :
\begin{itemize}
\item $t_\nu$ est un quantile de Student à $\nu=n_A-1$ degrés de liberté sur le court terme : l'écart-type repose sur des
dispersions estimées à partir d'une poignée de contributeurs, et une estimation gaussienne divisée par une échelle
indépendante de type $\chi^2_\nu$ suit une loi de Student ;
\item $s\ge1$ est une échelle calibrée sur un historique sain exactement comme $\kappa$, de sorte que les sous-jacents
sains déclenchent au plus $\alpha$ alertes d'erreur de marquage FO (\vLevelScaleEuStock, \vLevelScaleUsStock,
\vLevelScaleHtb{} et \vLevelScaleIndex{} pour les quatre catégories du banc d'essai).
\end{itemize}

\begin{example}[Une détection à la main]\label{ex-toy}
Quatre piliers $T=(0{,}25\,;0{,}5\,;1\,;2)$, un bruit indépendant d'écarts-types $(1\,;2\,;4\,;8)$\,pb de portage, un modèle
nul réduit au niveau de repo ($X=T$), un niveau vrai de $30$\,pb, un pic de $12$\,pb en $T_2$ et un bruit
$(0{,}4\,;-1{,}1\,;2{,}6\,;-3{,}0)$\,pb, d'où $y=(7{,}9\,;25{,}9\,;32{,}6\,;57{,}0)$\,pb. Le niveau ajusté vaut
$\hat\beta_0=\vToyBeta$\,pb et $Q=\vToyQ$ (le quantile à $99\,\%$ de $\chi^2_3$ vaut $11{,}3$). Les \vToyM{} gabarits donnent
\begin{center}\input{generated/\lang/tab_toy}\end{center}
La valeur critique familiale à $\alpha=2\,\%$ vaut $c=\vToyC$ (Šidák : \vToyCsidak). Le pic en $T_2$ l'emporte
($z=4{,}52$, amplitude $10{,}5\pm2{,}3$\,pb pour $12$ injectés) ; la marche à partir de $T_2$ obtient aussi un score élevé car
les deux gabarits sont corrélés. La poursuite adaptative retient le pic, après quoi tous les $|z|$ restants sont inférieurs
à $0{,}8$ et $Q$ tombe à $0{,}6$ : une seule explication suffit.
\end{example}

\subsection{Algorithmes}
\begin{algorithm}[ht]
\caption{Bootstrap du repo implicite}\label{alg-bootstrap}
\begin{algorithmic}[1]
\Require spot $S$, taux cumulé $R(\cdot)$, dividendes $(t_i,d_i)$, piliers $T_1<\dots<T_K$, forwards $F_1,\dots,F_K$
\For{$k=1,\dots,K$}
  \State construire $\Phi_k$ de \cref{eq-phik} à partir de $B_1,\dots,B_{k-1}$
  \State $x\gets R_k-\ln\bigl(F_k/(S-\hat D_k)\bigr)$ \Comment{dividendes de la période au repo précédent}
  \Repeat
    \State $g\gets\ln\Phi_k(x)-\ln F_k$ ; mettre à jour l'encadrement $[\ell,u]$ selon le signe de $g$
    \State $x\gets x-g\,\Phi_k(x)/\Phi_k'(x)$ ; si $x\notin(\ell,u)$, $x\gets(\ell+u)/2$
  \Until{le pas est inférieur à la tolérance}
  \State $B_k\gets x$
\EndFor
\State \Return $b_k=B_k/T_k$
\end{algorithmic}
\end{algorithm}

\begin{algorithm}[ht]
\caption{Détection structurelle}\label{alg-detect}
\begin{algorithmic}[1]
\Require écart de portage $y$, covariance $\Sigma$, échelle $\kappa$, modèle nul $X$, gabarits $\mathcal T$, niveau $\alpha$
\State $L\gets\operatorname{chol}(\kappa\Sigma)$ ; $P\gets$ projecteur qui retire $L^{-1}X$ ; $r\gets PL^{-1}y$ ; $Q\gets\|r\|^2$
\State pour chaque $t\in\mathcal T$ : $\tb\gets PL^{-1}t$, $z_t\gets\tb^\top r/\|\tb\|$
\State $c_\alpha\gets$ quantile $1-\alpha/2$ de $\max_t|Z_t|$, $Z\sim\N(0,C)$ \Comment{\cref{prop-family}}
\State alerte $\gets\max_t|z_t|>c_\alpha$ \textbf{ou} $Q>\chi^2_{K-2}(1-\alpha/2)$
\If{alerte}
  \State $\mathcal S\gets\emptyset$ ; \textbf{répéter} ajouter le gabarit de plus grand $|z|$ après projection de $[X,\mathcal S]$
  \Statex \hspace{\algorithmicindent} \textbf{tant que} ce $|z|$ dépasse $c_\alpha$
\EndIf
\State ajustement MCG conjoint de $y$ sur $[X,\mathcal S]$ ; amplitudes en unités naturelles ; portage réparé $y-\sum\hat a_tt$
\State alerte d'erreur de marquage FO si $|\hat\beta_0|>\max\bigl(10\text{\,pb},\;t_{\nu}\,s\operatorname{se}(\hat\beta_0)\bigr)$
\end{algorithmic}
\end{algorithm}

\begin{algorithm}[ht]
\caption{Calibration de l'échelle de variance}\label{alg-kappa}
\begin{algorithmic}[1]
\Require sous-jacents sains $i=1,\dots,N$ avec $(\max_t|z_{t,i}|,\,c_i,\,Q_i)$ calculés pour $\kappa=1$, cible $\alpha$
\State $\ell\gets\ln10^{-3}$, $u\gets\ln10^{3}$
\For{$60$ itérations}
  \State $m\gets(\ell+u)/2$ ; taux $\gets\frac1N\sum_i\one\bigl\{\max|z_i|e^{-m/2}>c_i\ \text{ou}\ Q_ie^{-m}>q_{1-\alpha/2}\bigr\}$
  \State \textbf{si} taux $\le\alpha$ \textbf{alors} $u\gets m$ \textbf{sinon} $\ell\gets m$
\EndFor
\State \Return $\kappa=e^{u}$
\end{algorithmic}
\end{algorithm}

\subsection{Exemples sur des sous-jacents synthétiques}
La \cref{fig-decomposition} montre le mécanisme dans l'espace de portage pour une erreur sur une période de dividendes,
et la \cref{fig-gallery} six détections dans l'espace de repo, avec les explications du détecteur au
\cref{tab-gallery}.

\begin{figure}[ht]
\centering\includegraphics[width=\linewidth]{figures/\lang/decomposition.pdf}
\caption{Vue en espace de portage de l'erreur sur une période de dividendes Totem de la \cref{fig-gallery}(b) : écart de
portage observé, ajustement nul avec sa bande de $\pm2$ écarts-types, et ajustement nul plus le gabarit de marche retenu.}
\label{fig-decomposition}
\end{figure}

\begin{figure}[p]
\centering\includegraphics[width=\linewidth]{figures/\lang/gallery.pdf}
\caption{Six actions européennes synthétiques, une par scénario, avec les réglages du détecteur du banc d'essai : repo
FO, repo implicite Totem, repo implicite après retrait de l'anomalie détectée, et bande de bruit par pilier. En (f), la
courbe Totem est saine et le repo FO est décalé : aucune anomalie n'est signalée et l'erreur de marquage FO est
publiée ; la bande est par pilier alors que le test d'erreur de marquage regroupe tous les piliers.}
\label{fig-gallery}
\end{figure}

\begin{table}[ht]
\centering
\caption{Première explication fournie par le détecteur pour chaque panneau de la \cref{fig-gallery}.}\label{tab-gallery}
\input{generated/\lang/tab_gallery}
\end{table}

\subsection{Les petits réseaux de neurones sont des filtres adaptés}
\begin{proposition}\label{prop-nn}
Supposons le modèle nul projeté et le bruit blanchi, de sorte que $r\sim\N(0,P)$ sous $H_0$ et $r\sim\N(a\tb,P)$ sous $H_1$,
avec un gabarit $t$ et une amplitude $a$ connus. Le logarithme du rapport des probabilités a posteriori vaut
\begin{equation}\label{eq-nn}
\ln\frac{P(H_1\mid r)}{P(H_0\mid r)}=a\,\tb^\top r-\tfrac12a^2\|\tb\|^2+\ln\frac{\pi_1}{\pi_0},
\end{equation}
une fonction affine du filtre adapté $\tb^\top r$. Un classifieur entraîné par entropie croisée sur suffisamment de données
converge vers cette probabilité a posteriori : dans ce cadre, un « petit réseau par type d'anomalie » ne peut au mieux
qu'apprendre le filtre adapté ; avec une amplitude et un signe inconnus, on retrouve la statistique du rapport de
vraisemblance $|z_t|$ du \cref{thm-glrt}.
\end{proposition}
\begin{proof}
Les deux densités gaussiennes ont la même covariance $P$, donc leur rapport logarithmique est affine en $r$ : les termes
quadratiques se compensent. L'entropie croisée est minimisée par la vraie probabilité a posteriori. Maximiser
\cref{eq-nn} en $a$ donne $(\tb^\top r)^2/(2\|\tb\|^2)$.
\end{proof}

\begin{insight}
Le détecteur est le détecteur optimal pour les mécanismes que l'on sait nommer, écrit plutôt qu'appris : son taux de
fausses alertes est fixé exactement, sa puissance est connue en forme fermée (\cref{cor-power}), et chaque alerte est
accompagnée d'une amplitude en unités naturelles. Les modèles appris trouvent leur place sur ce que les gabarits ne
couvrent pas (\cref{sec-alternatives}).
\end{insight}
